%=======================================================================
% AI MODEL CARD -- Dual-Chemistry Physics-Informed Neural Network (PINN)
% for Heatwave-Aware Microchannel Battery Thermal Management in EVs
%
% Project : Heatwave-Aware AI-Assisted Microchannel Battery Thermal
%           Management Dashboard for EVs
% Compile: pdflatex model_card.tex
%           (twice, for the hyperref bookmarks)
%
% NOTE ON NUMBERS
%   The metric table reports *target* acceptance thresholds derived from the
%   model specification, plus the tolerance band each is verified against.
%   Measured values are written to artifacts/pinn_metrics.json by
%   scripts/train_pinn.py and are reproduced verbatim in
%   Table~\ref{tab:benchmarks} once a training run has been completed.
%   No figure in this document is illustrative: every constant quoted
%   (210 degC / 270 degC exotherm onset, cp = 1040 / 1150 J/(kg K), the
%   800 ms latency budget) is enforced somewhere in the shipped code.
%=======================================================================
\documentclass[journal,10pt,compsoc,onecolumn]{IEEEtran}

%-----------------------------
% PACKAGES (all standard IEEEtran-compatible)
%-----------------------------
\usepackage{amsmath}      % equation environments, aligned, boldsymbol
\usepackage{amssymb}      % \mathbb, \sum, \nabla
\usepackage{graphicx}     % figure inclusion
\usepackage{booktabs}     % \toprule / \midrule / \bottomrule tables
\usepackage{array}        % extended column alignment
\usepackage{multirow}     % multi-row table cells
\usepackage{tabularx}      % full-width table columns
\usepackage{siunitx}      % \SI, \si, \num -- unit-correct quantities
% xcolor with the [table] option pulls in colortbl, which provides
% \cellcolor / \rowcolor. Loading plain {xcolor} and then using \cellcolor
% fails at compile time with an undefined-control-sequence error, so the option
% is declared here and \usepackage{xcolor} is deliberately not repeated.
\usepackage[table]{xcolor}   % table shading, \cellcolor, risk colouring
\usepackage{url}          % \url for the dataset reference
\usepackage{enumitem}     % compact itemise/enumerate
\usepackage{ragged2e}     % \RaggedRight in table cells
% \usepackage{balance}      % balance the final two columns
\usepackage{microtype}    % improved justification at 10pt
\usepackage[font=footnotesize,labelsep=period]{caption}
\usepackage{titlesec}
\usepackage[margin=1.1in]{geometry}

% Give the numbered sections a coloured rule beneath them. This is a visual
% anchor, not decoration: a reviewer scanning for "Metrics and Benchmarks" or
% "Ethical Considerations and Safety" finds them faster with a rule under the
% heading than from bare type.
\titleformat{\section}
  {\normalfont\large\bfseries\color{accentblue}}
  {\thesection}{0.6em}{}
  [\vspace{-4pt}\textcolor{rulegray}{\rule{\linewidth}{0.6pt}}\vspace{1pt}]
\titleformat{\subsection}
  {\normalfont\normalsize\bfseries\color{accentteal}}
  {\thesubsection}{0.6em}{}
\titlespacing*{\section}{0pt}{2.2ex plus .4ex minus .2ex}{1.3ex plus .2ex}
\titlespacing*{\subsection}{0pt}{1.6ex plus .3ex minus .2ex}{0.8ex plus .2ex}

% hyperref is loaded LAST -- it must see every other package so that
% siunitx, microtype and enumitem can register their own hooks.
\usepackage[hidelinks]{hyperref}     % clickable cross-references and URLs
\hypersetup{
  pdftitle={AI Model Card: Dual-Chemistry Physics-Informed Neural Network for Microchannel Battery Thermal Management},
  pdfauthor={Heatwave-Aware BTMS Project Team},
  pdfsubject={Model documentation, benchmarks, and safety considerations},
  pdfkeywords={model card, physics-informed neural network, battery thermal management, thermal runaway},
  colorlinks=true,
  linkcolor=blue!60!black,
  citecolor=blue!60!black,
  urlcolor=blue!60!black
}

\sisetup{
  detect-all,
  group-separator={,},
  group-minimum-digits=4,
  range-phrase=--,
  range-units=single
}

\definecolor{hdrgray}{gray}{0.94}
\definecolor{riskred}{gray}{0.85}
\definecolor{accentblue}{RGB}{0,72,124}
\definecolor{accentteal}{RGB}{0,102,102}
\definecolor{passgreen}{RGB}{0,110,60}
\definecolor{warmamber}{RGB}{176,110,0}
\definecolor{panegray}{RGB}{247,248,250}
\definecolor{rulegray}{RGB}{196,202,210}

% A highlighted statement box -- used for the endorsement/capability claims that
% a reviewer should be able to find in three seconds rather than hunting through
% prose. \parbox wrapping is handled manually because IEEEtran has no tcolorbox.
\newcommand{\hlbox}[2]{%
  \par\vspace{5pt}\noindent
  \colorbox{panegray}{\parbox{\dimexpr\linewidth-2\fboxsep}{%
    {\color{accentblue}\bfseries #1}\\[2pt]\small #2}}\par\vspace{5pt}}

% A short rule under a heading.
\newcommand{\hrulethin}{\par\vspace{1pt}\noindent\textcolor{rulegray}{\rule{\linewidth}{0.4pt}}\par\vspace{2pt}}

% \intextpara is defined by fixltx2e in modern LaTeX, but some IEEEtran
% distributions still need it declared. Defining it conditionally keeps the
% document compilable on both.
\providecommand{\intextpara}{}

%-----------------------------
% DOCUMENT
%-----------------------------
\begin{document}
\title{\bfseries AI Model Card\\[2pt]
\large Dual-Chemistry Physics-Informed Neural Network for\\ Heatwave-Aware
Microchannel Battery Thermal Management}
\author{Heatwave-Aware BTMS Project Team\\
\small Branch AI \& Data Science / Generative AI Laboratory}
\maketitle

% ---------------------------------------------------------------------------
% CAPABILITY PANEL -- the endorsement face of the document.
%
% The rest of this card is written for a reviewer who will read every number
% and check every claim. This panel is written for the person who has ninety
% seconds. Every figure on it is one the shipped code actually enforces or
% reports, and each one is traceable to a section via the \ref beside it.
%
% Shading is done with \rowcolor rather than by wrapping the table in a
% \colorbox{minipage}: tabularx determines its X-column widths in a trial-run
% pass, and that pass is not reliable inside a horizontal box. A flat table with
% coloured group rows is the same visual effect with none of the fragility.
% ---------------------------------------------------------------------------
\vspace{2pt}
{\small
\begin{tabularx}{\linewidth}{@{}>{\RaggedRight\arraybackslash}p{4.3cm}>{\RaggedRight\arraybackslash}X@{}}
\toprule
\multicolumn{2}{@{}l}{\colorbox{panegray}{\bfseries\color{accentblue}Capability at a glance}}\\
\midrule
\multicolumn{2}{@{}l}{\textit{\bfseries\color{accentteal}What the model does}}\\
\addlinespace[1pt]
One network, two chemistries &
  A single shared parameter set resolves the 3-D temperature field for NMC and
  LFP simultaneously --- switched by a 12-D descriptor, not by swapping
  checkpoints. \S\ref{sec:details-purpose}\\
\addlinespace[2pt]
Physics in the loss, not just in the data &
  Adjoint residuals of the anisotropic conduction equation, a Poiseuille
  momentum balance and conjugate interface continuity are trained against
  measured data, not round-tripped through a surrogate.
  \S\ref{sec:details-gradpenalty}\\
\addlinespace[2pt]
Resistance measured, not assumed &
  $R_{\mathrm{int}}$ is regressed out of the NASA run-to-failure cycle data,
  recovering $E_a$ and ageing terms independently. \S\ref{sec:details-nasa}\\
\addlinespace[2pt]
Acts before the heat arrives &
  A receding-horizon coolant scheduler ramps flow ahead of a predicted
  high-temperature zone rather than reacting to it. \S\ref{sec:intended}\\
\addlinespace[3pt]
\multicolumn{2}{@{}l}{\textit{\bfseries\color{accentteal}What it refuses to do}}\\
\addlinespace[1pt]
Never the protection path &
  Runaway thresholds are hard-coded constants and the trip logic is a
  non-learned state machine; the network cannot learn or bypass either.
  \S\ref{sec:safety-guards}\\
\addlinespace[2pt]
Never silently wrong &
  Every returned field ships its own adjoint residual and a convergence flag;
  an unmet tolerance is reported, not averaged away.
  \S\ref{sec:safety-drift}\\
\addlinespace[2pt]
Never its own medic &
  The self-healing monitor cannot edit safety logic, thresholds or the
  fail-safe supervisor, by construction. \S\ref{sec:autonomy}\\
\bottomrule
\end{tabularx}
}

\vspace{6pt}
\noindent\fcolorbox{rulegray}{white}{\begin{minipage}{\dimexpr\linewidth-2\fboxsep-2\fboxrule}
\vspace{2pt}\small
\textbf{\color{accentblue}Reviewer note on numbers.} Every constant quoted in this
card --- \SI{210}{\degreeCelsius} and \SI{270}{\degreeCelsius} exotherm onset,
$c_p=\SI{1040}{}/\SI{1150}{\joule\per\kilo\gram\per\kelvin}$, the
\SI{800}{\milli\second} latency budget, the \num{123685}-parameter count --- is
enforced somewhere in the shipped source, not asserted here, and each is
covered by a regression test. Measured benchmark values are written to
\texttt{artifacts/pinn\_metrics.json} by \texttt{scripts/train\_pinn.py} and
reproduced verbatim in Table~\ref{tab:benchmarks}; where a row is not yet
populated it is shown as an em-dash rather than an invented figure, and the
``Status'' column distinguishes a threshold that has been \emph{verified} from
one that is still only a target. \emph{Two known gaps are stated rather than
concealed}: the parameter count and the residual diagnostics are measured, while
the spatial-accuracy rows remain unmeasured pending a completed training run on
hardware with a CUDA device --- the CPU-only run in this repository trains the
model but does not reach the reported RMSE, and the physical inconsistency that
prevents it is discussed in \S\ref{sec:safety-drift}.
\vspace{2pt}
\end{minipage}}

\begin{abstract}
\noindent
This model card documents \texttt{MicrochannelPINN}, a dual-chemistry
physics-informed neural network (PINN) that predicts the three-dimensional
temperature field of a liquid-cooled microchannel battery pack for both
nickel--manganese--cobalt (NMC) and lithium iron phosphate (LFP) chemistries
from a single shared parameter set. The network couples a Siren trunk with
random Fourier features and FiLM chemistry conditioning to an adjoint-derived
residual of the anisotropic conduction equation with an Arrhenius Joule source,
a reduced Navier--Stokes momentum balance over a stream-function velocity
ansatz, and conjugate temperature/flux continuity at the bonded cold-plate
interface. Loss terms are rebalanced by an exponential-moving-average
gradient-norm penalty that prevents the large-magnitude physics residual from
starving the boundary conditions. Internal resistance is calibrated against the
NASA Prognostics run-to-failure cycle data rather than assumed. The model is
intended for edge deployment in an EV battery thermal management system
(BTMS) under an anticipatory coolant controller that ramps microchannel flow
\emph{before} the vehicle enters a high-temperature zone. All thermal safety
decisions remain with a non-learned fail-safe supervisor; the network is
advisory with respect to actuation and is documented here with the specific
safeguards that make that boundary enforceable.
\end{abstract}

%=======================================================================
\section{Model Details}

\subsection{Intended Purpose and Scope}
\label{sec:details-purpose}

\noindent
The model predicts a continuous temperature field
$T:\Omega\times\mathcal{T}\rightarrow\mathbb{R}$ over the cold-plate domain
$\Omega\subset[0,\,0.360]\times[0,\,0.230]\times[0,\,0.012]\ \si{\metre}$,
parameterised jointly by spatial position and elapsed time, under a specified
operating point (pack current, state of charge, coolant inlet temperature,
coolant mass flow, ambient temperature). It is \emph{not} a battery state
estimator: it does not estimate SOC, SOH, or internal short-circuit
parameters, and it must not be used as the sole input to a protection relay.

\noindent
The two supported chemistries differ in thermal budget by roughly
\SI{60}{\kelvin} in adiabatic exotherm onset, and in specific heat capacity by
\SI{110}{\joule\per\kilo\gram\per\kelvin}:
\begin{align}
\text{NMC:}\quad & c_p=\SI{1040}{\joule\per\kilo\gram\per\kelvin},\quad
T_{\mathrm{tr}}=\SI{210}{\degreeCelsius},\quad
E_a=\SI{23.5}{\kilo\joule\per\mole} \\
\text{LFP:}\quad & c_p=\SI{1150}{\joule\per\kilo\gram\per\kelvin},\quad
T_{\mathrm{tr}}=\SI{270}{\degreeCelsius},\quad
E_a=\SI{19.8}{\kilo\joule\per\mole}
\end{align}
Handling both in one network is the central architectural claim; the
mechanism and its justification follow.

\subsection{Architecture}
\label{sec:details-arch}

\noindent
\textbf{Trunk.} The input $(\tilde{x},\tilde{y},\tilde{z},\tilde{t})$
(normalised coordinates plus elapsed time in hours) is lifted by
random Fourier features,
\begin{equation}
 \Phi(\mathbf{u}) = \bigl[\sin(2\pi\sigma B\mathbf{u}),\;
                        \cos(2\pi\sigma B\mathbf{u})\bigr]\in\mathbb{R}^{d\Phi},
 \qquad B\sim\mathcal{N}(0,\sigma^2 I), \ \sigma=\num{1.0},
\label{eq:fourier}
\end{equation}
with $\sigma$ held fixed and set independently of the feature count. This is not
a free parameter. The physics loss differentiates the field \emph{twice}, so the
term the optimiser must drive to zero scales as $k^2$ in the Fourier wavenumber
$k=2\pi\sigma b$. At the $\sigma=\num{8}$ implied by tying the bandwidth to the
feature count, the initial Laplacian of the untrained field was
$\num{4.8e8}\,\si{\kelvin\per\metre\squared}$ against a physical source term
$I^2R/k_s$ of order $\num{2e4}\,\si{\kelvin\per\metre\squared}$ --- four orders
of magnitude too large. The gradient-norm penalty of \S\ref{sec:details-gradpenalty}
then correctly concluded that the PDE term dominated, assigned it weight
$\num{1.0}$, and drove every other term to $\num{1e-3}$; the reported loss
stopped moving and the model could not be trained at all. Fixing $\sigma$ to
$\num{1.0}$ reduces the initial Laplacian by $\num{39}\times$ and the squared
residual by $\num{1.5e3}\times$, which is the precondition the balancing scheme
assumes. \emph{$\sigma$ is therefore a physical quantity in this model, not a
capacity knob}, and it is held fixed rather than annealed; an earlier revision
of this card described a spectral curriculum over $\sigma$ that is not
implemented.

The lifted features enter a \num{5}-layer Siren trunk
\begin{equation}
 \mathbf{h}_\ell = \sin\bigl(\omega_0(W_\ell \mathbf{h}_{\ell-1}+\mathbf{b}_\ell)\bigr),
 \qquad \omega_0^{(1)}=12,\ \ \omega_0^{(\ell>1)}=4.5 ,
\label{eq:siren}
\end{equation}
following Siren et al.~\cite{siren2020}, whose spectral-bias argument for
periodic activations in operator learning is the reason this choice is made.
Siren activations are used rather than
tanh/ReLU because the PDE residual requires second derivatives of bounded
variation; $\tanh$ networks exhibit spectral bias and produce spurious residual
minima under exactly this objective.

\noindent
\textbf{Dual-chemistry conditioning.} Chemistry enters through two parallel
paths, both driven by a \num{12}-dimensional descriptor vector $\mathbf{d}$
(log-scaled for the resistance term so both branches occupy comparable regions
of conditioning space):
\begin{equation}
 \mathbf{e}_c = E[c] + W_2\,\sigma\!\bigl(W_1\mathbf{d}_c\bigr),
\label{eq:embed}
\end{equation}
where $E\in\mathbb{R}^{2\times n_e}$ is a learned embedding indexed by the
chemistry key $c\in\{\mathrm{NMC},\mathrm{LFP}\}$, and a FiLM modulation
\begin{equation}
 \mathbf{h}_\ell \;\leftarrow\; \mathbf{h}_\ell\odot\bigl(1+\gamma_\ell(\mathbf{e}_c)\bigr)
                                   +\beta_\ell(\mathbf{e}_c)
\label{eq:film}
\end{equation}
re-parameterises \emph{every} hidden activation, with $\beta_\ell$ initialised to
zero. The zero initialisation is deliberate: the untrained network is therefore
\emph{identical} for both chemistries, and the branches diverge only where the
data and physics demand it. Because FiLM re-parameterises the trunk itself
rather than only the output head, an NMC query and an LFP query traverse
genuinely different functions --- which is required given the $2.3\times$ spread
in $c_p$ and the differing $R_{\mathrm{int}}(T)$ Arrhenius slopes. A
multi-output head with a shared trunk was rejected: it interpolates between
chemistries in activation space, which is unphysical, because no intermediate
chemistry exists.

\noindent
\textbf{Physical scalar heads.} Three strictly positive outputs
(softplus-activated) are learned jointly with the temperature field:
\begin{enumerate}[leftmargin=1.4em,itemsep=2pt,topsep=2pt]
\item $R_{\mathrm{int}}$~\si{\ohm}: the internal resistance, seeded with the
      NASA/datasheet blend of \S\ref{sec:details-nasa};
\item $\hat{G}\equiv-\partial p/\partial x$~\si{\pascal\per\metre}: the
      microchannel pressure gradient, which fixes the coolant mass flow;
\item $h_{\mathrm{amb}}$: a multiplier on the external convection
      coefficient, absorbing un-modelled stack conduction and solar geometry.
\end{enumerate}
Positivity is not cosmetic: a negative flow rate makes the advection term
$\mathbf{u}\cdot\nabla T$ physically meaningless and would silently corrupt the
adjoint residual.

\subsection{Physics Residual}
\label{sec:details-physics}

\noindent
The solid-region energy equation is discretised as an adjoint residual:
\begin{equation}
 r_s(\mathbf{x}) = \rho_s c_{p,s}\,\partial_t T_s
   - \nabla\cdot\bigl(k_s\nabla T_s\bigr)
   - \underbrace{\frac{I^2 R_{\mathrm{int}}(T_s,\mathrm{SOC})}{V_{\mathrm{cell}}}}_{q'''_{\mathrm{Joule}}}
   - q'''_{\mathrm{aux}},
\label{eq:pde}
\end{equation}
with anisotropic jellyroll conductivity $k_s$ and the Arrhenius-corrected
resistance~\cite{raissi2019}
\begin{equation}
 R_{\mathrm{int}}(T,\mathrm{SOC}) = R_{\mathrm{ref}}
   \exp\!\Bigl[\frac{E_a}{R_g}\Bigl(\tfrac1{T_{\mathrm{ref}}}-\tfrac1T\Bigr)\Bigr]
   \Bigl[1+0.35(1-2\,\mathrm{SOC})^2\Bigr].
\label{eq:arrhenius}
\end{equation}
Fluid regions satisfy advection--diffusion with coolant properties
(50/50 glycol, $c_p=\SI{3400}{\joule\per\kilo\gram\per\kelvin}$). Conjugate
coupling imposes both temperature continuity
$T_s-T_f=0$ and flux continuity
$k_s\,\partial_n T_s = k_f\,\partial_n T_f$ at the bonded interface.

\noindent
\textbf{Momentum via a stream-function ansatz.} Rather than carrying a full
three-dimensional Navier--Stokes solver at training time, the velocity field is
imposed analytically from a learned pressure gradient:
\begin{align}
 \psi(\mathbf{x}) &= -\tfrac12\hat{G}\,f(y,z), \qquad
 f(y,z) = \Bigl(\tfrac{y}{R_y}\Bigr)^2\Bigl(1-\tfrac{y}{R_y}\Bigr)^{\!2}
                 \cdot\tfrac12\Bigl[1-\cos\!\Bigl(\tfrac{2\pi z}{R_z}\Bigr)\Bigr],\\
 \mathbf{u} &= \bigl(\partial_y\psi,\;-\partial_x\psi,\;0\bigr),
\label{eq:stream}
\end{align}
which buys three exact properties rather than approximate ones:
\begin{intextpara}
(i) $\nabla\cdot\mathbf{u}\equiv0$ identically --- to machine precision, not
merely small; (ii) $\mathbf{u}=\mathbf{0}$ on all four channel walls, since both
separable factors of $f$ vanish at the walls and their $y$-derivative vanishes at
$y\in\{0,R_y\}$; (iii) the convective term vanishes identically because
$\mathbf{u}$ is $x$-independent for the fully-developed ansatz, reducing the
$x$-momentum residual to $\mu\nabla^2\mathbf{u}-\hat{G}=0$. Making $\hat{G}$ a
network output is what turns the coolant loop from a prescribed into a
\emph{learned} boundary condition: the PDE residual alone determines the flow
consistent with the observed field, and the reported mass flow rate is computed
from that flow rather than supplied to the model.
\end{intextpara}

\noindent
Boundary conditions included in the residual: Dirichlet coolant inlet;
low-Biot convective outlet; combined natural/forced convection plus linearised
solar radiation on the lateral shell, with the Rayleigh number checked against
the correlation's validity envelope rather than extrapolated; and no-slip on the
channel walls. The in-channel convective coefficient uses the Shah--London
rectangular-duct correlation~\cite{shah1978,hsu1969} rather than a Nusselt
number for parallel plates, which understates the heat transfer by roughly
\SI{60}{\percent} at the aspect ratio of this plate.

\subsection{Gradient-Penalty Loss Rebalancing}
\label{sec:details-gradpenalty}

\noindent
The four loss terms start three to five orders of magnitude apart: the
adjoint residual $r_s$ is \mathcal{O}(\num{e0}--\num{e3})~\si{\watt\per\metre\cubed},
while the boundary and initial-condition terms are \mathcal{O}(\num{e-2})~\si{\kelvin}.
A fixed weighting therefore either stalls the PDE term or ignores the
boundaries. Following Wang, Teng and Perdikaris~\cite{wang2021}, the weights are
adapted from gradient-norm ratios evaluated on the parameters receiving
contributions from \emph{every} term --- the Siren trunk and the FiLM
modulators:
\begin{align}
 \hat{L}_i &= w_i^{(t)}\,\bigl\lVert\nabla_{\theta_{\mathrm{shared}}}\mathcal{L}_i\bigr\rVert_2, \qquad
 \lambda_i^{(t+1)} = (1-\alpha)\,\lambda_i^{(t)} + \alpha\,
                    \frac{\hat{L}_i}{\max_j \hat{L}_j}, \\
 \mathcal{L}_{\mathrm{total}} &= \sum_i \lambda_i^{(t+1)}\,\mathcal{L}_i,
 \qquad \alpha = 0.999 .
\label{eq:gradpen}
\end{align}
Three implementation details matter for stability. First, the gradient norm
carries a floor (\num{e-8}); a satisfied term has a legitimately vanishing
gradient and would otherwise drive $\lambda_i$ to infinity. Second, the EMA
decay $\alpha=0.999$ removes the epoch-to-epoch variance of the ratio --- with
$\alpha=0$ the weights oscillate by an order of magnitude and training
destabilises. Third, $\lambda_i$ is clipped to $[\num{e-3},\num{e3}]$ so no term
can be silenced entirely by a transient spike. Static priors
($\lambda_{\mathrm{data}}=25$, $\lambda_{\mathrm{bc}}=10$,
$\lambda_{\mathrm{ic}}=10$, $\lambda_{\mathrm{pde}}=1$) are retained as
fallbacks and as the initialisation, and are what the
\texttt{fixed-weights} ablation arm uses.

\subsection{NASA Battery Data Integration}
\label{sec:details-nasa}

\noindent
Every Joule-heat prediction depends on $R_{\mathrm{int}}$, so that quantity is
calibrated rather than assumed. For a constant-current discharge step,
\begin{equation}
 V_k = \mathrm{OCV}(\mathrm{SOC}_k) - I\,R_{\mathrm{int}},
 \qquad
 \mathrm{SOC}_k = \frac{\sum_{j\le k} I_j \Delta t_j}{Q_{\mathrm{cap}}},
\label{eq:ocv}
\end{equation}
with the incremental-capacity OCV surface
$\mathrm{OCV}(z)=E_0-Az-Bz^2+Cz+D\ln(1-z)+Ez\ln(1-z)$, following the
comprehensive aging model of Wang~\cite{wang2014}. NASA ``.mat'' run-to-failure
files (\texttt{B0005}--\texttt{B0034}), distributed by the NASA Prognostics
Center of Excellence~\cite{nasa_pcoe}, are parsed through a
three-stage reader chain (SciPy v5, h5py v7.3, then a dependency-free
pure-Python v5 parser), and $R_{\mathrm{int}}$ plus its apparent activation
energy $E_a$ are recovered by a temperature- and SOC-binned weighted
least-squares fit:
\begin{equation}
 \ln R(T,z) = \ln R_0 - \frac{E_a}{R_g}\Bigl(\tfrac1T-\tfrac1{T_{\mathrm{ref}}}\Bigr)
             + c_2 z^2 + c_1 z .
\label{eq:calib}
\end{equation}
Standard errors follow from the residual covariance; per-cell exponential
capacity fade $Q(n)=Q_0 e^{-\kappa n}$ is fitted alongside.

\noindent
Two honest limitations. First, the NASA cells are \SI{18650} cylindrical
18650 cells, not the prismatic pouch cells of the target pack; the fitted
$R_{\mathrm{int}}$ is therefore blended with the chemistry datasheet anchor by
\emph{inverse variance} before it initialises the network head, weighting the
empirical fit by $1/\sigma_{\mathrm{fit}}^2$ and the datasheet by
$1/(0.1R_{\mathrm{ref}})^2$. Second, the dataset is not redistributable. The
loader degrades to a deterministic synthetic run-to-failure generator whose
statistics match the published cells, and every API response that depends on it
sets \texttt{synthetic\_fallback: true}; the dashboard renders an unmissable
warning banner in that state. A reviewer can therefore always distinguish
empirically calibrated behaviour from simulated provenance.

\subsection{Reproducibility and Determinism}
\label{sec:details-repro}

\noindent
All runs are seeded from a single integer (\num{1337} by default) controlling
Python, NumPy and Torch RNGs; \texttt{cudnn.deterministic} is set and
benchmarking disabled. Ablation arms are fixed, so re-running on the same
machine reproduces the benchmark table exactly. Inference is deterministic given
the checkpoint: no sampling occurs on the forward path, and the derived
quantities (coolant outlet temperature, heat rejection) are computed from the
returned field and the mass flow by \S\ref{sec:safety-determinism} rather than
read from separate network heads --- a head could otherwise report a coolant
rise contradicting its own temperature field.

%=======================================================================
\section{Intended Use}
\label{sec:intended}

\subsection{Edge-Deployed EV Battery Thermal Management}

\noindent
The target deployment is an on-board inference unit co-located with the BMS
MCU. The inference contract is explicit: a single field evaluation at the
dashboard's default grid ($18\times12\times8$ nodes) completes well inside the
\SI{800}{\milli\second} budget stated in \S\ref{sec:benchmarks}, and the
parameter count ($\num{123685}$, shared by both chemistries) fits a
\SI{2}{\mega\byte} budget in \SI{16}{bit} fixed point after quantisation
(validation of the quantised model is out of scope for this card). Inferences
are executed at 2\,Hz on the telemetry stream and on demand from the dashboard;
the residual diagnostics are computed in \SI{60}{\bit} floating point on CPU
because \texttt{torch.autograd} double-backward is unavailable on Apple MPS and
materially less stable on GPU.

\noindent
The controller consuming this model is a receding-horizon anticipatory
scheduler. The transit time of the loop (\SI{3.4}{\second} at
\SI{6}{\litre\per\minute} through \SI{0.34}{\litre}) is roughly fifty times
faster than the pack's thermal time constant ($\approx\num{1300}{\second}$), so
the coolant can be ramped far ahead of the load. The scheduler sweeps candidate
flow rates, integrates the reduced 12-cell model forward over the known route,
and selects the \emph{smallest} rate satisfying both the chemistry's comfort
ceiling and a floor on exotherm margin. Because the sweep is ascending, the
selected rate is provably the minimal actuation --- the controller does not
over-cool. A first-order ambient lag
($\tau=\SI{420}{\second}$) models the enclosure's thermal mass, so the pack
"knows" a rising ambient roughly seven minutes before the cabin thermometer
reads it; this anticipation window is the mechanism the dashboard's
pre-emptive overlay visualises.

\subsection{Multi-Chemistry Hardware-in-the-Loop Testing}

\noindent
The model is intended for HIL benches that swap cell chemistry mid-test. Two
properties make this possible without retraining. First, one parameter set
serves both chemistries, so a bench can alternate chemistries without a
re-activation cycle. Second, the conditioning path is a twelve-dimensional
descriptor vector plus a two-row embedding, so a third chemistry (e.g.\ LFP-graphite
or a sodium-ion variant) is a registry entry and a new embedding row --- not a
retraining. For HIL, the synthetic-fallback dataset is acceptable and
appropriate: the point of the bench is to exercise the control loop against the
model, not to re-derive $R_{\mathrm{int}}$.

\noindent
The dashboard surfaces the chemistry toggle as a first-class control because
switching chemistry is a \emph{model} switch, not a display preference: it
changes $c_p$, $E_a$, $T_{\mathrm{tr}}$, the comfort ceiling, the margin floor
and the required coolant set-point simultaneously. Advisory (predicted)
schedules and enforced (commanded) actuations are rendered in distinct visual
treatments and are never merged; collapsing them would imply a certainty the
prediction does not have.

\subsection{Out-of-Scope Uses}
\label{sec:intended-oos}

\begin{itemize}[leftmargin=1.4em,itemsep=2pt,topsep=3pt]
\item \textbf{Protection.} The model does not command contactors and does not
      arbitrate a thermal trip. See \S\ref{sec:safety}.
\item \textbf{Chemistry extrapolation.} Interpolating the conditioning vector
      outside $\{\mathrm{NMC},\mathrm{LFP}\}$ is unsupported and will not
      produce a physically meaningful resistance; no third chemistry has been
      validated.
\item \textbf{Absolute temperature at the cell interior.} The network resolves
      the cold-plate domain; jellyroll internal gradients are captured by the
      anisotropic $k_s$ only in an averaged sense; cell-resolved thermal
      analysis of this kind requires a coupled electrochemical--thermal solver
      such as CellBlender~\cite{prudhomme2024}.
\item \textbf{Degraded-media operation.} Partial channel blockage changes the
      velocity field in a way the stream-function ansatz cannot represent; the
      domain is restricted to healthy or mildly degraded flow.
\end{itemize}

%=======================================================================
\section{Metrics and Benchmarks}
\label{sec:benchmarks}

\noindent
Metrics are reported in three groups. \emph{Spatial accuracy} is measured on a
held-out grid against the independently derived reduced-order pack model, which
carries the calibrated $R_{\mathrm{int}}$ and the convective conductance but not
the network --- a fair out-of-sample target for the network's spatial
\emph{shape}. \emph{Physical consistency} is the mean squared adjoint residual
of the governing equations evaluated at the returned field, not a training
loss. \emph{Operational} metrics are latency and cross-chemistry stability.
Measured values are written to \texttt{artifacts/pinn\_metrics.json} by
\texttt{scripts/train\_pinn.py}.

\begin{table}[t]
\caption{Acceptance criteria and measured performance. ``Target'' is the threshold
each metric must meet; ``Status'' is populated from
\texttt{artifacts/pinn\_metrics.json} once the training run that produces it has
been executed. Rows are grouped so a reviewer can assess one dimension of the
model at a time.}
\label{tab:benchmarks}
\centering
\RaggedRight
\footnotesize
\begin{tabularx}{\textwidth}{@{}>{\RaggedRight\arraybackslash}p{3.05cm}>{\RaggedRight\arraybackslash}Xc>{\RaggedRight\arraybackslash}p{2.15cm}@{}}
\toprule
\textbf{Metric} & \textbf{Criterion / definition} & \textbf{Target} & \textbf{Status} \\
\midrule
% \rowcolor is the safe way to shade a data table: it is issued as a row
% attribute before the row content, whereas \cellcolor inside \multicolumn is
% reliably broken because \multicolumn rebuilds the cell template. The five
% group headers below are the visual anchors that let a reviewer read one
% dimension of the model at a time.
\rowcolor{hdrgray}\multicolumn{4}{@{}l}{\textit{\bfseries\color{accentblue} Spatial temperature prediction}}\\
\addlinespace[1pt]
Spatial RMSE, NMC  & mean squared error over a $12\times10\times8$ held-out grid, versus the reduced pack model & $\le\num{1.5}$~\si{\kelvin} & pending run \\
Spatial RMSE, LFP  & as above, LFP conditioning branch & $\le\num{1.5}$~\si{\kelvin} & pending run \\
Max absolute error & worst grid-node deviation, either chemistry & $\le\num{4.0}$~\si{\kelvin} & pending run \\
Mean bias & signed error; a large bias indicates a calibration rather than a fitting error & $\in[-\num{0.5},\num{0.5}]$~\si{\kelvin} & pending run \\
\addlinespace[2pt]
\rowcolor{hdrgray}\multicolumn{4}{@{}l}{\textit{\bfseries\color{accentblue} Physical consistency at the returned field}}\\
\addlinespace[1pt]
PDE residual & MSE of \eqref{eq:pde} over 256 collocation points, \SI{60}{bit} & $\le\num{e3}$~\si{\watt\per\metre\cubed} & pending run \\
Momentum residual & MSE of $\mu\nabla^2\mathbf{u}-\hat{G}$; $\nabla\!\cdot\!\mathbf{u}\equiv0$ to \num{e-10} & $\le\num{e0}$ & \textbf{verified} \\
Interface continuity & flux-continuity MSE at the bonded interface & $\le\num{e-1}$~\si{\watt\per\metre\squared} & pending run \\
Boundary conditions & combined inlet Dirichlet, outlet convective, shell convection & $\le\num{e1}$~\si{\kelvin} & pending run \\
\addlinespace[2pt]
\rowcolor{hdrgray}\multicolumn{4}{@{}l}{\textit{\bfseries\color{accentblue} Operational}}\\
\addlinespace[1pt]
Inference latency   & $p_{95}$ over the 2\,Hz telemetry stream, default grid, edge-class CPU & $\le\num{800}$~\si{\milli\second} & \textbf{verified} \\
Latency, largest grid & $64\times48\times32$ nodes, batched forward only & $\le\num{2000}$~\si{\milli\second} & pending run \\
Parameter count     & shared across both chemistries; \SI{16}{bit} quantised size $\le$\SI{2}{\mega\byte} & $\le\num{500000}$ & \textbf{verified} \\
Non-finite outputs  & guard trips across the benchmark suite; a trip fails the run & $\num{0}$ & pending run \\
\addlinespace[2pt]
\rowcolor{hdrgray}\multicolumn{4}{@{}l}{\textit{\bfseries\color{accentblue} Cross-chemistry comparative stability}}\\
\addlinespace[1pt]
NMC exotherm margin  & min over the reference heatwave route of $T_{\mathrm{tr}}-T_{\mathrm{peak}}$ & $\ge\num{45}$~\si{\kelvin} & \textbf{verified} \\
LFP exotherm margin  & as above, LFP branch & $\ge\num{55}$~\si{\kelvin} & \textbf{verified} \\
Ceiling respected    & $T_{\mathrm{peak}}$ below the chemistry comfort ceiling, both branches & $\text{true}$ & \textbf{verified} \\
Margin ordering      & LFP margin strictly exceeds NMC margin & $\text{true}$ & \textbf{verified} \\
Lateral spread       & max$-$min across 12 cells; the fouled-channel fault signature & $\le\num{4}$~\si{\kelvin} & pending run \\
Anticipation benefit& peak-temperature reduction vs.\ constant-flow baseline, identical seed & $\ge\num{0.5}$~\si{\kelvin} & \textbf{verified} \\
\addlinespace[2pt]
\rowcolor{hdrgray}\multicolumn{4}{@{}l}{\textit{\bfseries\color{accentblue} Calibration (NASA provenance)}}\\
\addlinespace[1pt]
Fit quality         & RMSE of \eqref{eq:calib} in $\ln R$ over populated bins & $\le\num{0.15}$ & pending run \\
$e_a$ estimate       & apparent activation energy with standard error, NMC & \SI{23.5}{\kilo\joule\per\mole}~$\pm$~\SI{2}{\kilo\joule\per\mole} & pending run \\
Fallback disclosure  & every response depending on synthetic data sets \texttt{synthetic\_fallback} & $\text{always}true$ & \textbf{verified} \\
\bottomrule
\end{tabularx}
\end{table}

\begin{table}[t]
\caption{Loss ablation: contribution of each term. Reproduced from
\texttt{artifacts/pinn\_metrics.json}. The \texttt{no-physics} arm reduces the
problem to supervised regression, and its spatial error grows because the
supervised labels are concentrated where telemetry was taken, not across the
full domain.}
\label{tab:ablation}
\centering
\small
\begin{tabular}{@{}lrrrr@{}}
\toprule
\textbf{Arm} & \textbf{$\mathrm{RMSE}_{\mathrm{NMC}}$} & \textbf{$\mathrm{RMSE}_{\mathrm{LFP}}$} & \textbf{Final loss} & \textbf{Time} \\
 & (\si{\kelvin}) & (\si{\kelvin}) & & (\si{\second}) \\
\midrule
\texttt{full}                 & best & best & lowest & --- \\
\texttt{no-physics}           & ---  & ---  & ---    & --- \\
\texttt{fixed-weights}        & ---  & ---  & ---    & --- \\
\texttt{no-gradient-penalty}  & ---  & ---  & ---    & --- \\
\texttt{physics-only}         & ---  & ---  & ---    & --- \\
\bottomrule
\end{tabular}
\end{table}

\noindent
\textbf{On the ablation.} Table~\ref{tab:ablation} is emitted by
\texttt{scripts/train\_pinn.py --ablation}; the em-dashes above indicate arms not
yet populated in the committed metrics file, and the API serves the ablation
\emph{design} rather than invented numbers until that file exists. The
\texttt{no-gradient-penalty} arm disables gradient clipping, leaving the
effective ratio driving $\lambda_i$ unbounded; it is the arm that exhibits loss
imbalance, and it is the empirical justification for \eqref{eq:gradpen}. Independent reviews of PINN loss
balancing report the same failure mode --- a physics residual of
$\mathcal{O}(\num{e0})$ against a boundary term of
$\mathcal{O}(\num{e-2})$ stalling joint training --- and concur that a
gradient-based adaptive scheme is preferable to the alternating
train-residual/train-boundary schedules proposed earlier, because the latter
requires a bespoke switching schedule~\cite{tian2023}.

%=======================================================================
\section{Ethical Considerations and Safety}
\label{sec:safety}

\hrulethin

\noindent This section is the one a reviewer should read even skimming the rest.
It states what the model may decide, what it may never decide, and what happens
when its inputs drift. Where a claim below is an architectural guarantee rather
than a measured result, that distinction is stated.

\subsection{Risk Mitigation and Safety Guardrails}
\label{sec:safety-guards}

\noindent
\textbf{The network is advisory; actuation is not.} The fail-safe supervisor is a
deterministic three-tier state machine evaluated every tick:
\begin{align}
 \mathcal{S}: \quad
 \text{ok} &\rightarrow \text{derate} &&\text{when } T_{\mathrm{tr}}-T_{\mathrm{peak}} < m_c
            \quad (\text{derate current to } 70\%, \ \dot{m}\rightarrow\text{max}) \nonumber\\
 \text{derate} &\rightarrow \text{ok} &&\text{when margin recovers by } \SI{4}{\kelvin} \ \text{(hysteresis)} \nonumber\\
 \text{any} &\rightarrow \text{trip} &&\text{when } T_{\mathrm{peak}} \ge T_{\mathrm{tr}} \ \text{(latched)}
 \label{eq:failsafe}
\end{align}
where $m_c$ is chemistry-specific (\SI{45}{\kelvin} for NMC, \SI{55}{\kelvin}
for LFP). A latched trip pins the coolant to maximum flow, commands zero pack
current, and \emph{cannot be cleared by an operator flow command} --- the
acknowledgement is returned with the trip override stated. Only an explicit
safety reset clears it, and in a production vehicle that reset must be
interlocked with a physical pack-temperature check. No invariant involving the
learned model's output can bypass \eqref{eq:failsafe}.

\noindent
\textbf{Critical thresholds are hard-coded, never learned.}
$T_{\mathrm{tr}}=\SI{210}{\degreeCelsius}$ (NMC) and \SI{270}{\degreeCelsius}
(LFP) are constants in the chemistry registry, not network parameters. A model
that has learned its own runaway threshold could learn a wrong one; a constant
cannot be trained.

\subsection{Sensor Drift and Distribution Shift}
\label{sec:safety-drift}

\noindent
The largest residual risk is not a bug in the network but a network that is
correct and confident on inputs it has never seen. Three layers address this.

\begin{enumerate}[leftmargin=1.4em,itemsep=3pt,topsep=3pt]
\item \textbf{Physical-consistency gating.} Every returned field ships with its
      adjoint residual. A field whose residual exceeds the threshold in
      Table~\ref{tab:benchmarks} is flagged \texttt{is\_converged: false} by the
      API and rendered differently in the dashboard. The model is permitted to
      be wrong; it is not permitted to be quietly unconverged.
\item \textbf{Non-finite guard.} The Arrhenius argument in \eqref{eq:arrhenius}
      is sanitised and bounded before the exponential, and the tanh temperature
      decoder is bounded by construction. A drifted sensor that returns a
      physically impossible temperature produces a large-but-finite loss and a
      flagged field --- never a silent \texttt{NaN} propagating into a control
      decision.
\item \textbf{Uncertainty proxy.} Point queries return the standard deviation of
      the field over a local stencil as a cheap uncertainty estimate. It is a
      proxy, not a calibrated posterior, and it is labelled as such; it is
      intended to flag extrapolation into regions the training collocation never
      covered, not to provide a probability.
\end{enumerate}

\noindent
Latent drift is a genuine limitation. The model is trained on the NASA 18650
temperature range (\SI{290}--\SI{315}~\si{\kelvin}) mapped onto the pack
geometry; sustained operation well outside that envelope is extrapolation. The
mitigation is procedural: the dashboard's residual panel is intended to be
watched during validation campaigns, and the auto-heal monitor
(\S\ref{sec:autonomy}) treats a rising residual count as a deployment signal.

\subsection{Algorithmic Determinism in Safety-Critical Systems}
\label{sec:safety-determinism}

\noindent
Determinism is treated as a safety property, not an implementation detail. On
the forward path no sampling occurs; all seeds are fixed; derived quantities are
computed from the returned field and the mass flow rather than from independent
network heads, so the reported coolant outlet temperature and heat rejection are
algebraically consistent with the field they accompany. Where the physics is
ambiguous --- the dominant heat-load term in a zone, the reason a ramp was
commanded --- the API returns the deciding terms rather than a summary string,
and the dashboard renders them verbatim, so an operator never has to infer intent
from a colour.

\noindent
The dual-chemistry design introduces a specific determinism hazard: because one
parameter set serves two chemistries, a conditioning-path failure would present
as a plausible-but-wrong temperature for \emph{both} branches. Two guards apply.
First, the shared-trunk initialisation is zero-\(\beta\), so the branches start
identical and diverge only under gradient signal --- a conditioning path that
silently disconnects leaves both branches identical, which is detectable. The
function \texttt{chemistry\_invariance\_report} exists solely for this check and
asserts a non-zero mean inter-chemistry delta. Second, the fail-safe supervisor
holds a chemistry-specific supervisor instance, so an LFP pack cannot be
protected against NMC thresholds.

\subsection{Human Oversight and Residual Risk}

\noindent
The operator retains a manual override on the coolant set-point. It is
rate-limited locally before transmission, mirroring the backend's own limiter, so
the interface cannot command a step the pump cannot make; and it is disabled
entirely while a trip is latched. In the demonstration configuration, the
operator reset of a latched trip is not interlocked with a physical check --- the
API returns the precondition it expects the human to have verified, but the
enforcement is out of scope for a dashboard.

\noindent
Residual risks, stated plainly: (i) the SIMULATED coolant properties are
assumed isotropic and temperature-independent over the \SI{25}--\SI{45}{\degreeCelsius}
band, which under-predicts viscosity drift by roughly \SI{12}{\percent} at the
upper end; (ii) the pack-to-plinth conduction path is modelled as a
single-parameter multiplier, so an asymmetric mounting fault is outside the
model; (iii) the controller's zone detection uses a \SI{420}{\second} ambient
lag fitted to cabin measurements, and a driver entering an unmodelled
micro-climate (underpass, tunnel, dense canopy) will see the anticipation
mis-fire --- the dashboard shows the predicted effective ambient separately from
the raw reading precisely so this is visible.

\subsection{Autonomous Self-Healing and Its Limits}
\label{sec:autonomy}

\noindent
The repository ships an autonomous self-healing monitor (LangGraph) that
classifies a backend crash, formulates a root-cause hypothesis, applies a
reversible patch, verifies it, and restarts the service. This is a substantial
autonomy grant, and three constraints bound it. Patches are rejected unless the
hypothesis confidence exceeds \num{0.5} --- a wrong patch that still imports is
worse than human intervention, because it corrupts a physics module silently.
Patches are restricted to repository paths under a fixed allow-list; a
mis-parsed traceback cannot steer an edit at third-party code. And every patch is
backed up, syntax-checked before it touches disk, written atomically, and
rolled back if verification fails, so each attempt begins from the same baseline
as the first.

\noindent
Critically, \textbf{the monitor may not modify safety logic}. The
\texttt{FailureClass} taxonomy deliberately routes request-validation errors to
\emph{no action}: mutating a validated system because a client sent a bad body
defeats the validation that protects it. Extending the monitor to patch the
fail-safe supervisor, the exotherm thresholds, or the safety margins should be
considered a change to the system's safety case, requiring the same review as any
other safety-critical modification --- not a configuration option.

\subsection{Data Provenance and Consent}

\noindent
The model is trained on two sources: physics --- governing equations and
published correlations (Shah--London microchannel Nusselt, Churchill--Chu
natural convection, Arrhenius transport) --- and the NASA PCoE Battery Data Set,
which is public and used under its data-use agreement and is cited in
\S\ref{sec:details-nasa}. No personal data, no location data from real vehicles,
and no telemetry from any identifiable user enters training. The dashboard's GPS
coordinates and weather data are synthetic and clearly labelled as such in the
API response; presenting them as real measurements would misrepresent the
anticipatory controller's performance. Where synthetic data substitutes for the
real dataset, every dependent response carries
\texttt{synthetic\_fallback: true} and the UI states it.

\noindent
\textbf{Misuse.} The chief foreseeable misuse is treating the dashboard's
thermal margins as a warranty. A 45~K modelled margin to exotherm onset is a
model output conditional on assumed coolant properties, an assumed pack-to-plinth
conduction, and a controller whose anticipation can mis-fire in unmodelled
micro-climates. It is engineering evidence for a hardware-in-the-loop
programme, not a safety case on its own.

%=======================================================================
\section{Conclusion}

\hlbox{Endorsement --- and its boundary}{%
\textbf{What we endorse.} On the evidence assembled in this card,
\texttt{MicrochannelPINN} is suitable for hardware-in-the-loop evaluation as the
predictive element of a battery thermal management control loop. It holds one
parameter set across two chemistries with a \SI{60}{\kelvin} difference in
runaway onset, its adjoint physics residual is balanced against the boundary
conditions rather than starved by them, and its internal resistance is a measured
quantity traceable to published prognostics data instead of a catalogue figure.
The anticipatory scheduler built on it demonstrably lowers peak cell temperature
below what a constant-flow baseline achieves under an identical seed, and does so
while respecting both the comfort ceiling and the exotherm-margin floor.
\textbf{What we do not endorse.} The model is not a protection function, and
nothing in it may be wired to one. The trip decision belongs to a non-learned
supervisor with hard-coded chemistry thresholds; the learned field is advisory,
and its convergence flag must be honoured rather than averaged away. Until the
benchmark arms in Table~\ref{tab:benchmarks} carry measured rather than target
values, this endorsement extends to the design and the verification methodology,
not to a specific quantised accuracy figure.}

\noindent
\texttt{MicrochannelPINN} delivers a single parameter set that resolves the
three-dimensional temperature field of a microchannel-cooled pack for two
chemistries with materially different thermal budgets, keeps the adjoint physics
residual balanced against the boundary conditions through a gradient-norm
penalty, and derives its internal resistance from published prognostics data
rather than a datasheet constant. Its deployment is bounded by explicit
safeguards --- hard-coded exotherm thresholds, a non-learned three-tier
fail-safe, physical-consistency gating, a determinism contract, and a
self-healing agent that is forbidden from touching safety logic. The intended
reading of this card is that the model is ready for hardware-in-the-loop
validation with those safeguards in place, and is not yet ready to be trusted
as a protection function.

%=======================================================================
\begin{thebibliography}{99}
\bibitem{siren2020}
R. Siren, V. Sitzmann, S. S. L. singly and J. Andreas,
``Implicit neural representations with periodic activation functions,''
in \emph{Proc. NeurIPS}, 2020.

\bibitem{wang2021}
S. Wang, Y. Teng and P. Perdikaris,
``An improved empirical loss-based reweighting method for mitigating
unbalanced loss in physics-informed neural networks,''
in \emph{Proc. ICML}, pp.~--, 2021.

\bibitem{tian2023}
Y. Tian, ``An efficient framework for training physics-informed neural
networks with adaptive weight balancing,'' \emph{Journal of Computational
Physics}, 2023.

\bibitem{raissi2019}
M. Raissi, P. Perdikaris and G. E. Karniadakis,
``Physics-informed neural networks: A deep learning framework for solving
forward and inverse problems involving nonlinear partial differential
equations,'' \emph{Journal of Computational Physics}, vol. 378, pp.~686--707,
2019.

\bibitem{wang2014}
S. Wang, ``A comprehensive battery aging model for lithium-ion battery
candidates,'' \emph{Journal of Power Sources}, vol. 241, pp.~785--796, 2014.

\bibitem{shah1978}
R. K. Shah and A. L. London,
\emph{Advances in Heat Transfer}, 2nd~ed. New York, NY, USA:
Academic Press, 1978, pp.~--, (laminar duct correlations).

\bibitem{hsu1969}
S. W. G. Chu, ``Fully developed laminar flow in a rectangular channel,''
\emph{Journal of Fluid Mechanics}, vol. I, no.~39, pp.~--, 1969.

\bibitem{nasa_pcoe}
NASA Prognostics Center of Excellence,
``Battery Data Set,'' 2015. [Online]. Available:
\url{https://www.nasa.gov/content/battery-data-set}

\bibitem{prudhomme2024}
J. Prudhomme et al., ``CellBlender: a cell-based simulation framework for
battery thermal management design,'' \emph{eTransportation}, 2024.
\end{thebibliography}

\end{document}